\chapter{Écrire en \LaTeX{} avec l'assistant}

Un assistant écrit vite du \LaTeX{} correct : il connaît les paquets, les environnements,
les macros d'une classe de revue. Il écrit aussi, par défaut, une prose reconnaissable entre
mille. Ce chapitre décrit la chaîne qu'on a suivie pour les articles du projet : un premier
jet par l'assistant, la compilation, une passe d'humanisation, puis la lecture de l'auteur.

\section{Le premier jet}

On donne à l'assistant le plan, les sources et les contraintes de la revue : longueur,
style de citation, classe, résumé. Il écrit les fichiers \fichier{.tex} et les compile
lui-même.

\begin{terminal}
# compiler sans installation de TeX ; les paquets se téléchargent au besoin
tectonic -X compile article.tex
\end{terminal}

\begin{regle}
L'assistant ne s'arrête pas à « ça compile » : il relit la sortie, corrige les lignes qui
débordent dans la marge et les références non résolues, et regarde les pages qu'il a
modifiées.
\end{regle}

Les nombres de l'article viennent de scripts, et les figures sont dessinées depuis les
données (chapitre~\ref{ch:publier}). L'assistant les relance quand les données changent ;
il n'invente pas un chiffre pour boucher un trou, et le dit quand il manque une source.

\section{Une passe d'humanisation}

La prose d'un modèle de langage a des tics : adjectifs gonflés, triades, phrases-slogans en
fin de paragraphe, tirets partout, connecteurs mécaniques. En français s'y ajoutent des
anglicismes (« adresser un problème ») et un vocabulaire administratif. Un relecteur les
remarque, et certaines revues s'en méfient.

On a écrit pour cela deux procédures, à la manière du chapitre précédent :
\texttt{humaniseur-fr} pour le français et \texttt{humanizer} pour l'anglais. Chacune
recense les tics, dit par quoi les remplacer, et impose quelques règles de prudence :

\begin{itemize}
  \item une seule passe ; une deuxième aplatit le texte au lieu de l'améliorer ;
  \item ne jamais abaisser le registre : un article reste soutenu ;
  \item ne toucher qu'à la prose, jamais aux formules, aux commandes, aux macros ni aux
        références ;
  \item ne rien ajouter qui ne soit dans la source ou dans les données : ni fait, ni
        chiffre, ni citation.
\end{itemize}

\begin{terminal}
# dans la conversation avec l'assistant
/humaniseur-fr docs/article/sections/*.tex   (article en français)
/humanizer docs/paper/sections/*.tex         (article en anglais)
\end{terminal}

Pour un article en anglais, on rédige directement en anglais et l'on passe
\texttt{humanizer}. Traduire un texte français déjà humanisé reproduit d'autres tics,
ceux de la traduction ; mieux vaut une passe sur le texte final, dans sa langue.

\begin{tipblock}
La passe d'humanisation trouve souvent des affirmations trop fortes en même temps que des
tics de style : « démontre » là où l'article cité « conjecture », « premier » là où
l'on n'a rien vérifié. On les corrige sur le fond, pas seulement sur la forme.
\end{tipblock}

\section{La lecture de l'auteur}

Rien ne remplace la dernière lecture, crayon en main. L'auteur y vérifie que chaque phrase
dit ce qu'il pense, que chaque référence dit ce qu'on lui fait dire, et qu'il saurait
défendre le texte devant un rapporteur sans l'assistant à côté. C'est aussi là qu'il réécrit
dans sa propre voix ce qu'il publiera sous son nom, comme l'exigent certaines revues et la
plupart des sites d'entraide.
